\exercisesection{深度}

\begin{enumerate}
\def\labelenumi{\arabic{enumi})}
\item
  设 A 为环，J 为 A 的理想，M 为 A-模。证明：为使 \(\operatorname{prof}_{\mathrm{A}}(J;M)\geqslant2\) 成立，必要且充分的是，由 J 到 A 中的典范单射所导出的映射 \(M\to\operatorname{Hom}_{\mathrm{A}}(J,M)\) 是一个同构。
\item
  设 A 为一个高为 1 的赋值环，且不是 Noether 的，m 为其极大理想，n 为 A 的一个主理想，异于 (0) 且异于 A。证明我们有 \(\operatorname{prof}_{\mathrm{A}}(\mathfrak{n};\mathrm{A})=1\) 以及 \(\operatorname{prof}_{\mathrm{A}}(\mathfrak{m};\mathrm{A})\geqslant2\) ，尽管 \(\mathrm{V}(\mathfrak{m})=\mathrm{V}(\mathfrak{n})\) （注意 \(\operatorname{Ext}_{\mathrm{A}}^{1}(\mathrm{A}/\mathfrak{m},\mathrm{A}/\mathfrak{n})\) 同构于 \(\operatorname{Hom}_{\mathrm{A}}(\mathrm{A}/\mathfrak{m},\mathrm{A})\) 。）
\item
  设 A 为局部环 \(k[[X,Y]]\) ，并设 M 为诸 A-模 \(\mathrm{A}/\mathfrak{a}\) 的直和，其中 \(\mathfrak{a}\) 取遍 A 的非零主理想组成的集合。证明有 \(\mathrm{prof}_{\mathrm{A}}(\mathrm{M}) = 1\) ，但在 \(\mathfrak{m}_{\mathrm{A}}\) 中不存在 M-正则的元素。
\item
  设 A 为环，J 为 A 的一个有限生成理想，M 为 A-模，P 为平坦 A-模。证明有 \(\operatorname{prof}_{\mathrm{A}}(\mathrm{J};\mathrm{P}\otimes_{\mathrm{A}}\mathrm{M})\geqslant\operatorname{prof}_{\mathrm{A}}(\mathrm{J};\mathrm{M})\) ，且当 P 忠实平坦时等号成立。
\item
  设 \(\mathrm{A}\) 为 Noether 环，M 为一个非零的有限生成 A-模，J 为 \(\mathrm{A}\) 的理想，\(\mathbf{x} = (x_1, \ldots, x_n)\) 为 J 的一个生成组。证明有 \(\mathrm{H}^i(\mathbf{x}, \mathrm{M}) \neq 0\) （对 \(\operatorname{prof}_{\mathrm{A}}(\mathrm{J}; \mathrm{M}) \leqslant i \leqslant n\) ）（化为局部情形，并利用 A, X, 第 157 页，推论 2）。
\item
  设 A 为 Noether 局部环，并设 P 为由平坦 A-模组成的一个长度有限为 \(\ell\) 的复形，使得 \(\operatorname{Supp}(\mathrm{H}(\mathrm{P})) = \{\mathfrak{m}_{\mathrm{A}}\}\) 。证明每个使得 \(\kappa_{\mathrm{A}} \otimes_{\mathrm{A}} \mathrm{M} \neq 0\) 的 A-模 M 的深度都 \(\leqslant \ell\) （把命题 3 应用于复形 \(\mathrm{P} \otimes_{\mathrm{A}} \mathrm{M}\) ，注意到这个复形不是正合的，并利用习题 4）。
\item
  设 A 为局部环，M 为有限生成 A-模，F 为 \(\operatorname{Spec}(A)\) 的一个闭子集。证明不等式 \(\operatorname{prof}_{\mathbf{F}}(\mathbf{M}) \geqslant \operatorname{prof}(\mathbf{M}) - \dim(\mathbf{F})\) 。
\item
  设 A 为 Noether 环，并设 M 为一个有限生成 A-模。我们称 M 满足性质 \((S_{k})\) ，如果对 A 的每个素理想 p 有 \(\operatorname{prof}_{A_{\mathfrak{p}}}(M_{\mathfrak{p}}) \geqslant \inf(k, \dim_{A_{\mathfrak{p}}}(M_{\mathfrak{p}}))\) 。
\end{enumerate}

\begin{enumerate}
\def\labelenumi{\alph{enumi})}
\item
  每个模都满足 \((\mathrm{S}_{k})\) （对 \(k \leqslant 0\) ）；说一个模满足 \((\mathrm{S}_{1})\) ，意味着它没有嵌入的相伴素理想。
\item
  设 \(0 \to M' \to M \to M'' \to 0\) 为有限生成 A-模的一个短正合列，并设 \(k\) 为整数。我们假设 \(M\) 满足 \((S_k)\) ，\(M''\) 满足 \((S_{k-1})\) ，且当 \(k \geqslant 2\) 时有 \(\operatorname{Ass}(M'') \subset \operatorname{Ass}(M)\) 。证明 \(M'\) 满足 \((S_k)\) （利用第 7 节的推论 2）。
\item
  c) 设 \((L,p)\) 为 M 的一个由有限生成自由 A-模组成的分解。若 A 满足性质 \((S_{k})\) ，则 A-模 \(B_{i}(L)\) （对 \(i \geqslant 0\) ）满足 \((S_{h})\) ，其中 \(h = \inf(k,i+1)\) 。 d) 若 A 满足 \((S_{2})\) ，则每个有限生成 A-模的对偶都满足 \((S_{2})\) （应用 c）。 e) 设 J 为 A 的一个理想，由一个 M-正则序列 \((x_{1},\ldots,x_{r})\) 生成。若 M 满足 \((S_{k})\) ，则 A-模 M/JM 满足 \((S_{k-r})\) 。
\end{enumerate}

\(f)\) 设 B 为 Noether A-代数，N 为一个有限生成 B-模且是忠实平坦的 A-模，并设 \(k\) 为整数。若 B-模 \(\mathrm{M} \otimes_{\mathrm{A}} \mathrm{N}\) 满足性质 \((\mathrm{S}_k)\) ，则 A-模 M 亦满足。若 M 满足 \((\mathrm{S}_k)\) ，且 \((\kappa(\mathfrak{p}) \otimes_{\mathrm{A}} \mathrm{B})\) -模 \(\kappa(\mathfrak{p}) \otimes_{\mathrm{A}} \mathrm{N}\) 满足 \((\mathrm{S}_k)\) （对每个 \(\mathfrak{p} \in \operatorname{Supp}(\mathrm{M})\) ），则 B-模 \(\mathrm{M} \otimes_{\mathrm{A}} \mathrm{N}\) 满足 \((\mathrm{S}_k)\) 。

\begin{enumerate}
\def\labelenumi{\arabic{enumi})}
\setcounter{enumi}{8}
\item
  给出一个局部环 A 和一个有限生成 A-模 M 的例子，使得 \(\mathfrak{m}_{\mathrm{A}}M \neq M\) 且 \(\operatorname{prof}_{A}(M) = +\infty\) （取 A 为一个以无限族未定元的多项式环的局部化）。
\item
  设 A 为 Noether 环，M 为有限生成 A-模，J 为 A 的理想，\((x_{1},\ldots,x_{r})\) 为 J 中元素组成的一个 M-正则序列，且 \(r < \operatorname{prof}_{\mathrm{A}}(\mathrm{J};\mathrm{M})\) 。设 \(q_{1}\ldots,q_{n}\) 为 A 的不含 J 的理想，除可能的两个外均为素理想，并设 \(J_{0}\) 为 J 的一个子集，它对加法和乘法稳定且生成理想 J。证明存在 \(J_{0}\) 的一个元素 x，它不属于任何 \(q_{i}\) ，且使序列 \((x_{1},\ldots,x_{r},x)\) 是 M-正则的（利用 II, § 1, 第 1 节，命题 2 的推论 2）。
\item
  设 \(\mathrm{A}\) 为环，M 为 \(\mathrm{A}\) -模，且 \(x_{1},\ldots ,x_{r}\) 为 \(\mathrm{A}\) 的元素。
\end{enumerate}

\begin{enumerate}
\def\labelenumi{\alph{enumi})}
\item
  对每个满足 \(1 \leqslant p \leqslant r\) 的整数 p，令 \(M_{p} = M/(x_{2}M + \ldots + x_{p}M)\) 。证明：若序列 \((x_{1}, \ldots, x_{r})\) 是 M-正则的，则序列 \((x_{1}, x_{p+1}, x_{p+2}, \ldots, x_{r})\) 是 \(M_{p}\) -正则的（对 p 用归纳法论证）。
\item
  为使序列 \((x_{1}, x_{3})\) 与 \((x_{2}, x_{3})\) 都是 M-正则的，必要且充分的是序列 \((x_{1}x_{2}, x_{3})\) 是 M-正则的。
\item
  设 \(x_p' \in \mathrm{A}\) （ \(1 \leqslant p \leqslant r\) ）。为使序列 \((x_1, \ldots, x_p, \ldots, x_r)\) 与 \((x_1, \ldots, x_p', \ldots, x_r)\) 都是 M-正则的，必要且充分的是序列 \((x_1, \ldots, x_p x_p', \ldots, x_r)\) 是 M-正则的（利用 a) 与 b)）。
\end{enumerate}

\(d\) 设 \(n_1,\ldots ,n_r\) 为 \(\geqslant 1\) 的整数。为使序列 \((x_{1}^{n_{1}},\dots,x_{r}^{n_{r}})\) 是 M-正则的，必要且充分的是序列 \((x_{1},\ldots ,x_{r})\) 是 M-正则的。

\begin{enumerate}
\def\labelenumi{\arabic{enumi})}
\setcounter{enumi}{11}
\tightlist
\item
  设 A 为环，M 为 A-模。
\end{enumerate}

\begin{enumerate}
\def\labelenumi{\alph{enumi})}
\item
  设 \(P \in A[X]\) 。证明：若 M{[}X{]}中以 P 为比的乘法映射的核不为零，则它含有 M 的一个非零元素（设 Q 为 M{[}X{]}中使 PQ = 0 的一个次数最小的元素；若 \(P = a_{0}X^{p} + \ldots + a_{p}\) ，\(Q = m_{0}X^{q} + \ldots + m_{q}\) ，注意 \(a_{0}Q = 0\) ，然后用归纳法注意对所有 i 有 \(a_{i}Q = 0\) ）。
\item
  设 J 为 A 的一个有限生成理想。证明关系 \(\operatorname{prof}_{\mathrm{A}}(\mathrm{J};\mathrm{M}) > 0\) 等价于存在一个 M{[}X{]} -正则的元素（在 JA{[}X{]}中）。
\item
  设 I 与 J 为 A 的有限生成理想。证明等式 \(\mathrm{prof}_{\mathrm{A}}(\mathrm{IJ};\mathrm{M}) = \min (\mathrm{prof}_{\mathrm{A}}(\mathrm{I};\mathrm{M}),\mathrm{prof}_{\mathrm{A}}(\mathrm{J};\mathrm{M}))\) （对 \(\mathrm{prof}_{\mathrm{A}}(\mathrm{IJ};\mathrm{M})\) 用归纳法推理，并利用 \(b)\) ）。
\end{enumerate}

设 A 为环，J 为 A 的理想，M 为 A-模。对每个整数 \(n \geqslant 0\) ，用 \(\operatorname{prof}_{n}(\mathbf{J}; \mathbf{M})\) 表示诸序列长度的上确界（在 \(\overline{\mathbf{N}}\) 中），这些序列由 \(\mathrm{M}[\mathrm{X}_{1}, \ldots, \mathrm{X}_{n}]\) -正则的、属于 \(\mathrm{JA}[\mathrm{X}_{1}, \ldots, \mathrm{X}_{n}]\) 的元素组成。

\begin{enumerate}
\def\labelenumi{\alph{enumi})}
\item
  证明序列 \(\left(\operatorname{prof}_{n}(\mathrm{J};\mathrm{M})\right)_{n\geqslant0}\) 是递增的；用 \(\operatorname{prof}_{\infty}(\mathrm{J};\mathrm{M})\) 表示它的极限（在 \(\overline{\mathbf{N}}\) 中）。证明我们有 \(\operatorname{prof}_{\infty}(\mathrm{J};\mathrm{M})\leqslant\operatorname{prof}_{\mathrm{A}}(\mathrm{J};\mathrm{M})\) 。
\item
  当理想 J 是有限生成的时，证明等式 \(\operatorname{prof}_{\infty}(\mathrm{J};\mathrm{M})=\operatorname{prof}_{\mathrm{A}}(\mathrm{J};\mathrm{M})\) （按第 4 节定理 2 的证明方式推理，并利用习题 12。）
\item
  在一般情形，证明 \(\operatorname{prof}_{\infty}(J;M)\) 是诸数 \(\operatorname{prof}_{\mathrm{A}}(J';M)\) 的上确界，其中 \(J'\) 取遍含于 J 的有限生成理想组成的集合。
\item
  设 \(J'\) 为 A 的一个理想，使得 \(V(J') = V(J)\) 。证明我们有 \(\text{prof}_{\infty}(J'; M) = \text{prof}_{\infty}(J; M)\) 。
\item
  设 A 为一个高为 1 的赋值环，且不是 Noether 的，并设 m 为其极大理想。证明我们有 \(\operatorname{prof}_{\infty}(\mathfrak{m};\mathrm{A})=1\) 和 \(\operatorname{prof}_{\mathrm{A}}(\mathfrak{m};\mathrm{A})\geqslant2\) （参见习题 2）。
\end{enumerate}

\begin{enumerate}
\def\labelenumi{\arabic{enumi})}
\setcounter{enumi}{13}
\tightlist
\item
  设 A 为一个完备 Noether 局部环。
\end{enumerate}

\begin{enumerate}
\def\labelenumi{\alph{enumi})}
\item
  证明：若 A 的一个理想含于 A 的一个由素理想组成的序列之并，则它含于其中的一个（注意拓扑空间 A 是一个 Baire 空间，参见 TG, IX, 第 55 页，定理 1）。
\item
  设 M 为一个具有可数生成族的 A-模。证明定理 2 的结论仍然成立（注意集合 Ass(M) 是可数的，并由 a) 推出：若 \(\mathrm{prof}_{\mathrm{A}}(\mathrm{J};\mathrm{M}) > 0\) ，则存在 J 的一个元素 \(x\) ，使得乘法映射 \(x_{\mathrm{M}}\) 是单射）。
\end{enumerate}

\begin{enumerate}
\def\labelenumi{\arabic{enumi})}
\setcounter{enumi}{14}
\tightlist
\item
  设 A 为 Noether 局部环，并设 M 为一个有限生成 A-模。证明不等式
\end{enumerate}

\[
\operatorname{prof} (A) \leqslant \operatorname{grade} (M) + \dim (M) \leqslant \dim (A).
\]

（对整数 \(g = \text{grade}(M)\) 用归纳法证明第一个不等式，利用第 7 节命题 13 的推论 2 处理 \(g = 0\) 的情形，然后考虑 Ann(M) 的一个 A-正则元素。至于第二个不等式，注意有 grade(M) ≤ prof(A \(_{p}\) )（对每个 p ∈ Supp(M)）。）

\begin{enumerate}
\def\labelenumi{\arabic{enumi})}
\setcounter{enumi}{15}
\tightlist
\item
  设 A 为 Noether 局部环。对每个有限生成 A-模 M，用 \(t_{\mathrm{A}}(\mathrm{M})\) 或简称 \(t(\mathrm{M})\) 表示 \(\kappa_{A}\) -向量空间 \(\operatorname{Ext}_{\mathrm{A}}^{p}(\kappa_{\mathrm{A}},\mathrm{M})\) 的维数，其中 \(p=\operatorname{prof}(\mathrm{M})\) 。
\end{enumerate}

\begin{enumerate}
\def\labelenumi{\alph{enumi})}
\item
  设 \((x_{1},\ldots,x_{r})\) 为 \(m_{A}\) 中对 M 完全割的元素组成的一个序列，它生成一个理想 J。我们有 \(t(\mathrm{M})=t(\mathrm{M}/\mathrm{JM})\) 。
\item
  设 \(\rho : A \to B\) 为 Noether 局部环之间的一个局部同态；设 M 为一个有限生成 A-模，并设 N 为一个在 A 上平坦的有限生成 B-模。证明公式 \(t_{B}(M \otimes_{A} N) = t_{A}(M)t_{B}(\kappa_{A} \otimes_{A} N)\) （利用 \(a\) 化为 \(\mathrm{prof}_{\mathrm{B}}(M \otimes_{A} N) = \mathrm{prof}_{\mathrm{A}}(M) = \mathrm{prof}_{\mathrm{B}}(\kappa_{\mathrm{A}} \otimes_{\mathrm{A}} N) = 0\) 的情形）。
\end{enumerate}

\begin{enumerate}
\def\labelenumi{\arabic{enumi})}
\setcounter{enumi}{16}
\tightlist
\item
  设 A 为 Noether 环，M 为有限生成 A-模。证明 M 是自反的，当且仅当下列两个条件被满足：
\end{enumerate}

\begin{enumerate}
\def\labelenumi{(\roman{enumi})}
\item
  对每个 \(\mathfrak{p} \in \operatorname{Spec}(A)\) （满足 \(\operatorname{prof}(A_{\mathfrak{p}}) \leqslant 1\) ），\(A_{\mathfrak{p}}\) -模 \(M_{\mathfrak{p}}\) 是自反的；
\item
  对每个 \(\mathfrak{p} \in \operatorname{Spec}(\mathrm{A})\) （满足 \(\operatorname{prof}(A_{\mathfrak{p}}) \geqslant 2\) ），有 \(\operatorname{prof}_{\mathrm{A}_{\mathfrak{p}}} (M_{\mathfrak{p}}) \geqslant 2\) 。
\end{enumerate}

（在条件 (i) 与 (ii) 之下，证明典范同态 \(M \to M^{**}\) 的核、然后其余核，都没有相伴素理想。）

\begin{enumerate}
\def\labelenumi{\arabic{enumi})}
\setcounter{enumi}{17}
\tightlist
\item
  设 A 为 Noether 环，J 为 A 的理想，M 与 N 为有限生成 A-模，且 \(\operatorname{prof}_{\mathrm{A}}(J;N)\geqslant2\) 。
\end{enumerate}

\begin{enumerate}
\def\labelenumi{\alph{enumi})}
\item
  证明不等式 \(\mathrm{prof}_{\mathrm{A}}(\mathrm{J};\mathrm{Hom}_{\mathrm{A}}(\mathrm{M},\mathrm{N}))\geqslant 2\) （考虑 J 中的一个 N-正则序列 \((x,y)\) ）。
\item
  假设 \(\mathrm{prof}_{\mathrm{A}}(\mathrm{J};\mathrm{Hom}_{\mathrm{A}}(\mathrm{M},\mathrm{N}))\geqslant 3\) 。证明我们有 \(\mathrm{prof}_{\mathrm{A}}(\mathrm{J};\mathrm{Ext}_{\mathrm{A}}^{1}(\mathrm{M},\mathrm{N}))\geqslant 1\) （考虑与 J 的一个 N-正则元素相伴随的扩张模的正合列）。
\end{enumerate}

